\hypertarget{appendix-k-deep-dive-the-memory-layout-of-a-c-class}{%
\section{Appendix K: Deep Dive: The Memory Layout of a C++
Class}\label{appendix-k-deep-dive-the-memory-layout-of-a-c-class}}

To become a C++ God, you must be able to ``see'' the memory. You should
be able to look at a class definition and sketch out its byte-by-byte
layout in your head.

Let's dissect a complex \textbf{Multiple Inheritance} hierarchy and see
how the compiler (GCC/Clang) arranges it in RAM.

\hypertarget{the-lab-rat-a-multiple-inheritance-hierarchy}{%
\subsection{The Lab Rat: A Multiple Inheritance
Hierarchy}\label{the-lab-rat-a-multiple-inheritance-hierarchy}}

\begin{lstlisting}[language={C++}]
class A {
    int a;
public:
    virtual void f() { std::cout << "A::f"; }
};

class B {
    int b;
public:
    virtual void g() { std::cout << "B::g"; }
};

class C : public A, public B {
    int c;
public:
    virtual void f() override { std::cout << "C::f"; } // Overrides A::f
    virtual void h() { std::cout << "C::h"; }          // New virtual function
};
\end{lstlisting}

\begin{center}\rule{0.5\linewidth}{0.5pt}\end{center}

\hypertarget{visualizing-class-c-in-memory}{%
\subsection{1. Visualizing Class C in
Memory}\label{visualizing-class-c-in-memory}}

Assuming a 64-bit system (where pointers are 8 bytes and
\passthrough{\lstinline!int!} is 4 bytes).

\hypertarget{the-object-layout-of-c}{%
\subsubsection{\texorpdfstring{The Object Layout of
\texttt{C}}{The Object Layout of C}}\label{the-object-layout-of-c}}

\begin{lstlisting}
[ Offset ] [ Size ] [ Content ]
--------------------------------------------------
[ 0      ] [ 8    ] [ vptr_A ]  --> Points to vtable for C (A-part)
[ 8      ] [ 4    ] [ int a  ]  -- From Class A
[ 12     ] [ 4    ] [ padding]  -- Alignment to 8-byte boundary
[ 16     ] [ 8    ] [ vptr_B ]  --> Points to vtable for C (B-part)
[ 24     ] [ 4    ] [ int b  ]  -- From Class B
[ 28     ] [ 4    ] [ int c  ]  -- From Class C
--------------------------------------------------
Total Size: 32 bytes
\end{lstlisting}

\hypertarget{why-the-padding}{%
\subsubsection{Why the Padding?}\label{why-the-padding}}

The CPU likes to read 8-byte chunks (on a 64-bit machine). If an 8-byte
pointer (\passthrough{\lstinline!vptr\_B!}) started at an odd address
like 12, the CPU would have to do two memory reads to get one pointer.
The compiler adds \textbf{padding} at offset 12 to ensure
\passthrough{\lstinline!vptr\_B!} starts at offset 16 (a multiple of 8).

\begin{center}\rule{0.5\linewidth}{0.5pt}\end{center}

\hypertarget{the-virtual-tables-vtables}{%
\subsection{2. The Virtual Tables
(Vtables)}\label{the-virtual-tables-vtables}}

Since \passthrough{\lstinline!C!} inherits from both
\passthrough{\lstinline!A!} and \passthrough{\lstinline!B!}, it actually
has \textbf{two} vtable pointers.

\hypertarget{vtable-for-c-primary---a-part}{%
\subsubsection{Vtable for C (Primary - A
part)}\label{vtable-for-c-primary---a-part}}

This vtable is used when you have an
\passthrough{\lstinline!A* ptr = new C();!}.

\begin{lstlisting}
[ Index ] [ Content ]
--------------------------------------------------
[ 0     ] [ C::f()  ]  -- Overridden
[ 1     ] [ C::h()  ]  -- New function in C is appended here
\end{lstlisting}

\hypertarget{vtable-for-c-secondary---b-part}{%
\subsubsection{Vtable for C (Secondary - B
part)}\label{vtable-for-c-secondary---b-part}}

This vtable is used when you have a
\passthrough{\lstinline!B* ptr = new C();!}.

\begin{lstlisting}
[ Index ] [ Content ]
--------------------------------------------------
[ 0     ] [ B::g()  ]  -- Not overridden
[ 1     ] [ thunk to C::f() ] -- Magic!
\end{lstlisting}

\hypertarget{what-is-a-thunk}{%
\subsubsection{What is a ``Thunk''?}\label{what-is-a-thunk}}

When you call \passthrough{\lstinline!ptr->f()!} through a
\passthrough{\lstinline!B*!}, the pointer is pointing to the
\emph{middle} of the object (offset 16). But
\passthrough{\lstinline!C::f()!} expects the
\passthrough{\lstinline!this!} pointer to point to the \emph{start} of
the object (offset 0). A \textbf{thunk} is a tiny piece of assembly that
subtracts 16 from the \passthrough{\lstinline!this!} pointer before
jumping to the real \passthrough{\lstinline!C::f()!}.

\begin{center}\rule{0.5\linewidth}{0.5pt}\end{center}

\hypertarget{data-alignment-rules-the-golden-ratio}{%
\subsection{3. Data Alignment Rules (The Golden
Ratio)}\label{data-alignment-rules-the-golden-ratio}}

\begin{enumerate}
\def\labelenumi{\arabic{enumi}.}
\tightlist
\item
  \textbf{Fundamental Alignment}: Every type has an alignment
  requirement. \passthrough{\lstinline!char!} is 1,
  \passthrough{\lstinline!short!} is 2, \passthrough{\lstinline!int!} is
  4, \passthrough{\lstinline!double/pointers!} are 8.
\item
  \textbf{Member Alignment}: A member must start at an offset that is a
  multiple of its alignment.
\item
  \textbf{Class Alignment}: The total size of the class must be a
  multiple of its \emph{largest} member's alignment.
\end{enumerate}

\hypertarget{example-of-wasteful-layout}{%
\subsubsection{Example of Wasteful
Layout:}\label{example-of-wasteful-layout}}

\begin{lstlisting}[language={C++}]
class Waste {
    char a;   // 1 byte
    double b; // 8 bytes
    char c;   // 1 byte
};
// Layout: [a] [7 bytes padding] [bbbbbbbb] [c] [7 bytes padding]
// Total: 24 bytes
\end{lstlisting}

\hypertarget{optimized-layout}{%
\subsubsection{Optimized Layout:}\label{optimized-layout}}

\begin{lstlisting}[language={C++}]
class Lean {
    double b; // 8 bytes
    char a;   // 1 byte
    char c;   // 1 byte
    // 6 bytes padding
};
// Total: 16 bytes (Saved 8 bytes!)
\end{lstlisting}

\textbf{Godhood Tip}: Always declare your members from largest to
smallest to minimize padding waste.

\begin{center}\rule{0.5\linewidth}{0.5pt}\end{center}

\hypertarget{how-to-inspect-this-yourself}{%
\subsection{4. How to Inspect This
Yourself}\label{how-to-inspect-this-yourself}}

Want to see the truth? Use the compiler's secret flags:

\textbf{For Clang:}

\begin{lstlisting}[language=bash]
clang++ -Xclang -fdump-record-layouts -c my_file.cpp
\end{lstlisting}

\textbf{For GCC:}

\begin{lstlisting}[language=bash]
g++ -fdump-lang-class my_file.cpp
\end{lstlisting}

This will output the exact byte offsets the compiler is using. Don't
take my word for itverify it with the machine!
